\section{Conclusion}
\label{sec:conclusion}

Multi-hop reasoning is fundamentally an information-acquisition problem rather than merely a generation problem. A system must determine what is already known, identify what remains unresolved, decide what information should be sought next, and combine the resulting evidence into a supported answer or report. Longer outputs, larger context windows, or more tool calls do not by themselves constitute multi-hop capability. The defining property is dependency: later retrieval and reasoning steps become meaningful because of information acquired earlier.

This survey organized the field into four functions. \emph{Initial retrieval} establishes an entry point from the original query. \emph{Follow-up retrieval} converts intermediate evidence into new information needs and search actions. \emph{Evidence integration} preserves, organizes, and combines information gathered across hops. \emph{Feedback learning} uses evaluations and prior trajectories to improve future retrieval and reasoning. The first three functions explain how a system completes the current task; the fourth explains how repeated task execution can produce durable improvement. This separation makes it possible to compare multi-hop QA pipelines, graph retrieval, adaptive RAG, agentic search, deep research, memory, and skill-based systems without treating them as unrelated research areas or reducing them to a single architectural template.

The evaluation implications are equally important. Final-answer accuracy cannot reveal whether the system found the required evidence, followed a valid search path, preserved key relations under a context budget, actually used the cited sources, stopped efficiently, or learned a reusable strategy. Progress therefore requires process-level evaluation of retrieval depth, hop completion, trajectory quality, evidence membership and ordering, grounding, cost, and cross-task improvement. Benchmarks and reporting practices should make these intermediate objects observable whenever a paper claims to improve them.

The next generation of multi-hop systems will be better characterized as adaptive information-seeking systems than as a retriever attached to a language model. Such systems will need to search, reason, verify, manage context, preserve provenance, and reuse experience under changing information environments. The central research objective is not simply to make models produce more elaborate reasoning traces, but to make them reliably acquire the right information, recognize unresolved uncertainty, integrate evidence without losing support, and improve their behavior from previous attempts.
